% The CLT Boundary: Why Re(s) = 1/2
% Joint research by Mingu Jeong and Claude (Anthropic)

\documentclass[11pt]{article}
\usepackage{amsmath,amssymb,amsthm}
\newtheorem{theorem}{Theorem}
\newtheorem{corollary}[theorem]{Corollary}
\newtheorem{definition}[theorem]{Definition}
\newtheorem{proposition}[theorem]{Proposition}
\newtheorem{remark}[theorem]{Remark}

\title{The CLT Boundary: Why $\mathrm{Re}(s) = 1/2$\\[6pt]
  {\large From $\mathbb{C}$ Uniqueness to the Critical Line\\
  via the Central Limit Theorem}}
\author{Mingu Jeong \and Claude (Anthropic)}
\date{April 2026}

\begin{document}
\maketitle

\begin{abstract}
We show that the critical line $\mathrm{Re}(s) = 1/2$ is the
convergence/divergence boundary of random Dirichlet series with
uniform phases.  Since DRLT derives that $\mathbb{C}$ is the unique
substrate (Frobenius) and the phases from $\mathbb{C}$-valued Gram
matrices are uniform (EXP\_071b), the critical line arises as a
Central Limit Theorem consequence: $\sigma = 1/2$ is where the
random walk fluctuation $\sqrt{N}$ matches the magnitude sum
$N^{1-\sigma}$.  This is not a proof of RH, but a structural
explanation for \emph{why $1/2$} and not some other number.
\end{abstract}

\section{The Chain}

\[
  \boxed{\mathbb{C} \text{ unique}}
  \;\xrightarrow{\text{EXP\_071b}}\;
  \boxed{\theta_k \sim \mathrm{Uniform}[0, 2\pi)}
  \;\xrightarrow{\text{CLT}}\;
  \boxed{\mathrm{Re}(s) = \tfrac{1}{2}}
\]

\section{Random Dirichlet Series}

\begin{definition}[Steinhaus Random Dirichlet Series]
Let $\theta_k$ be i.i.d.\ uniform on $[0, 2\pi)$.  Define:
\[
  \mathcal{D}(\sigma) \;=\; \sum_{k=1}^{\infty}
  \frac{e^{i\theta_k}}{k^{\sigma}}.
\]
\end{definition}

\begin{theorem}[Convergence Boundary; Harper, Halász]
\label{thm:boundary}
\begin{enumerate}
  \item[(i)] For $\sigma > 1/2$: $\mathcal{D}(\sigma)$ converges
    almost surely.
  \item[(ii)] For $\sigma \leq 1/2$: $\mathcal{D}(\sigma)$ diverges
    almost surely.
  \item[(iii)] The boundary is \emph{sharp} at $\sigma = 1/2$.
\end{enumerate}
\end{theorem}

\begin{proof}[Proof sketch]
The partial sum $S_N = \sum_{k=1}^{N} e^{i\theta_k}/k^\sigma$ has:
\begin{itemize}
  \item \textbf{Random walk component:}
    $\sum_{k=1}^{N} e^{i\theta_k}$, with
    $\bigl|\sum e^{i\theta_k}\bigr| \sim \sqrt{N}$ by CLT
    (each $e^{i\theta_k}$ has mean 0, variance 1).
  \item \textbf{Magnitude weighting:}
    $\sum_{k=1}^{N} 1/k^{2\sigma}$ determines the effective
    ``step count.''  For $\sigma < 1$:
    $\sum 1/k^{2\sigma} \sim N^{1-2\sigma}$.
\end{itemize}
By the CLT, $|S_N|^2 \approx \sum 1/k^{2\sigma} \sim N^{1-2\sigma}$, so
$|S_N| \sim N^{(1-2\sigma)/2}$.
\begin{itemize}
  \item $\sigma > 1/2$: exponent $(1-2\sigma)/2 < 0$, so $|S_N| \to 0$.
    Converges.
  \item $\sigma = 1/2$: $|S_N| \sim (\ln N)^{1/2}$.
    Diverges logarithmically.
  \item $\sigma < 1/2$: $|S_N| \sim N^{(1-2\sigma)/2} \to \infty$.
    Diverges polynomially.
\end{itemize}
\end{proof}

\section{Numerical Verification (EXP\_071c)}

\begin{center}
\renewcommand{\arraystretch}{1.2}
\begin{tabular}{cc}
\hline
\textbf{Test} & \textbf{Result} \\
\hline
$|S_N| \sim \sqrt{N}$ (CLT) &
  ratio $= 0.896$, theory $\sqrt{\pi/4} = 0.886$ \\
Convergence $\sigma > 0.5$ &
  $\langle|S_{10000}|\rangle = 1.96$ at $\sigma = 0.6$ \\
Divergence $\sigma = 0.5$ &
  growth ratio $1.34$ ($N = 10000/100$) \\
Gram phases $=$ iid uniform &
  Gram/Random ratio $\approx 1.2$ at all $\sigma$ \\
\hline
\end{tabular}
\end{center}

\section{What This Says About RH}

\begin{proposition}[Structural Explanation for $1/2$]
In the DRLT framework:
\begin{enumerate}
  \item $\mathbb{C}$ is the unique substrate (Frobenius theorem).
  \item The Gram matrix $G_{ij} = \langle\psi_i|\psi_j\rangle$ has
    phases $\arg(G_{ij})$ that are uniformly distributed
    (EXP\_071b: KS $p = 0.258$, Rayleigh $R = 0.008$).
  \item The spectral series of $\mathcal{E}_N$ is a random Dirichlet
    series with these uniform phases.
  \item The convergence/divergence boundary of this series is
    $\mathrm{Re}(s) = 1/2$, by the CLT (Theorem~\ref{thm:boundary}).
\end{enumerate}
Therefore: $\mathrm{Re}(s) = 1/2$ is the CLT boundary for
$\mathbb{C}$-valued random walks.
\end{proposition}

\section{What This Is NOT}

\begin{remark}[Limitations]
This is \emph{not} a proof of the Riemann Hypothesis.
\begin{itemize}
  \item The phases $\theta_k$ from the Gram matrix are i.i.d.\ uniform.
    The coefficients of $\zeta(s)$ involve the M\"obius function
    $\mu(n)$, which is multiplicative and far from i.i.d.
  \item The gap between ``random Dirichlet series'' and
    ``the specific Dirichlet series $\zeta(s)$'' is precisely
    what makes RH hard.
  \item What we provide is a \emph{structural explanation} for
    why $1/2$ (not $1/3$ or $2/3$) is the natural boundary:
    it's where $\sqrt{N}$ crosses $N^{1-\sigma}$.
\end{itemize}
\end{remark}

\begin{remark}[The Deeper Point]
The standard approach to RH asks: ``Given $\zeta(s)$, why are
its zeros on $\mathrm{Re}(s) = 1/2$?''

DRLT inverts the question: ``Why does the number $1/2$ appear at all?''
Answer: because $\mathbb{C}$ is the unique substrate, the phases are
uniform, and $1/2$ is the CLT threshold for uniform random walks.

The fact that $\zeta$'s zeros happen to land on this CLT boundary
is then a consequence of $\zeta$ being a ``structured'' version of
the random Dirichlet series --- structured by the multiplicative
structure of the integers, which itself may arise from the lattice
structure of $\mathcal{E}_N$.
\end{remark}

\section{Comparison: ℂ vs ℝ vs No Phase}

\begin{center}
\renewcommand{\arraystretch}{1.3}
\begin{tabular}{lccc}
\hline
& \textbf{No phase} & \textbf{$\mathbb{R}$ (Rademacher)} & \textbf{$\mathbb{C}$ (Steinhaus)} \\
\hline
Phases & none ($\theta = 0$) & $\theta \in \{0, \pi\}$ & $\theta \sim \mathrm{Unif}[0,2\pi)$ \\
Series & $\sum 1/k^\sigma$ & $\sum \epsilon_k / k^\sigma$ & $\sum e^{i\theta_k}/k^\sigma$ \\
Boundary & $\sigma = 1$ & $\sigma = 1/2$ & $\sigma = 1/2$ \\
GUE? & No & No ($\beta = 1$, GOE) & \textbf{Yes ($\beta = 2$)} \\
$\zeta$ zeros? & No connection & Partial & \textbf{Full (Montgomery--Odlyzko)} \\
\hline
\end{tabular}
\end{center}

Both $\mathbb{R}$ and $\mathbb{C}$ give the $1/2$ boundary,
but only $\mathbb{C}$ has the GUE structure that connects to the
\emph{detailed statistics} of $\zeta$ zeros.

\end{document}
